\section{Site-level PTM-omics extractors for DIG Gene Set Extractors}
\label{sec:ptm_site_level}

\subsection{Scope}
This note defines the site-level PTM-omics extractors currently implemented in
\texttt{geneset\_extractors}:

\begin{enumerate}
\item \texttt{ptm\_site\_diff}: convert a precomputed site-level contrast table
into a gene-weight vector.
\item \texttt{ptm\_site\_matrix}: convert a site-by-sample abundance matrix plus
sample metadata into one or more within-study site contrasts, then reuse the
same downstream scorer.
\end{enumerate}

The intended primary use case is phosphoproteomics from CPTAC/PDC-style public
site tables, but the scoring stack also applies to other PTM types when the
input contains a stable site identifier or enough fields to reconstruct one.

The central object is a gene-level program that captures differential PTM
occupancy at sites assigned to a gene, optionally adjusted for matched total
protein abundance and optionally downweighted for highly ubiquitous sites.

\subsection{Workflow W1: public CDAP/PDC report standardization}
\label{sec:ptm_public_prepare}

Public CPTAC/PDC phosphoproteomics often starts from CDAP-style isobaric text
reports rather than from a repo-standard site matrix. The implemented workflow
\texttt{ptm\_prepare\_public} performs only \emph{structural} harmonization:

\begin{enumerate}
\item detect the phosphosite report, optional matched proteome report, and
optional sample-design table;
\item parse each public phosphosite label into either a single resolved
canonical site key or a deterministic site-group key;
\item preserve the original public strings
(\texttt{raw\_site\_label}, \texttt{raw\_peptide}) together with a
\texttt{site\_parser\_status};
\item standardize sample columns and metadata into repo-native
\texttt{ptm\_matrix.tsv}, \texttt{sample\_metadata.tsv}, and optional
\texttt{protein\_matrix.tsv};
\item emit provenance/QC sidecars that can feed directly into
\texttt{ptm\_prepare\_reference\_bundle}.
\end{enumerate}

This workflow intentionally does \emph{not} perform alias resolution,
ubiquity weighting, protein-adjusted scoring, duplicate collapse, or gene
aggregation. Those remain downstream responsibilities of
\texttt{ptm\_site\_diff}, \texttt{ptm\_site\_matrix}, and the optional PTM
reference bundle.

\paragraph{Canonicalized public-site outputs.}
For a public row $r$, the helper attempts to emit either
\[
\texttt{site\_id} = (\mathrm{protein\_accession}, \mathrm{residue},
\mathrm{position}, \mathrm{ptm\_type})
\]
when a single site is losslessly parsable, or
\[
\texttt{site\_group\_id} = (\mathrm{protein\_accession},
\mathrm{site\_group\_label}, \mathrm{ptm\_type})
\]
for unresolved multi-site groups. If neither can be reconstructed losslessly,
the helper emits a deterministic fallback site-group key derived from the raw
public label and marks the row with
\texttt{site\_parser\_status = unparsed\_fallback}.

\subsection{What the implemented extractor captures}
\begin{itemize}
\item \texttt{ptm\_site\_diff} captures a compact gene program from an already
computed site-level contrast. With signed scores, positive weights indicate
genes whose measured PTM sites are relatively more modified in the focal
condition and negative weights indicate the opposite. With protein adjustment,
the program is interpreted as PTM-specific regulation beyond bulk protein
abundance changes.
\end{itemize}

\subsection{Notation}
\begin{itemize}
\item Raw site-level rows are indexed by $i \in \{1,\ldots,M\}$.
\item Canonical sites (or site-groups if the source cannot be localized to a
single residue) are indexed by $k \in \{1,\ldots,K\}$.
\item Genes are indexed by $g \in \{1,\ldots,G\}$.
\item $a(i) = k$ is the alias-resolution map from raw row $i$ to canonical site
key $k$.
\item $g(k)$ is the resolved gene assignment for canonical site $k$.
\item $d_i$ is a site-level effect size, typically a log fold-change.
\item $t_i$ is a site-level test statistic, when available.
\item $p_i$ is a site-level $p$-value or adjusted $p$-value, when available.
\item $\ell_i \in [0,1]$ is a site-localization confidence, when available.
\item $m_i$ is a peptide or evidence count, when available.
\item $a_{g(k)}^{\mathrm{prot}}$ is a matched total-protein effect for the gene
owning site $k$, when available.
\item $u_k$ is a reference ubiquity weight for site $k$, derived from an
optional phosphosite reference bundle.
\item $\psi(\cdot)$ is the user-selected signed or nonnegative score transform.
\item $S_g$ is the final gene score before selection and normalization.
\item $w_g$ is the selected and normalized gene weight written to
\texttt{geneset.tsv}.
\end{itemize}

\subsection{Canonical site identity and aliasing}
Public phosphoproteomics tables often use heterogeneous feature labels:
gene-symbol plus residue-position, UniProt accession plus residue-position,
study-specific site-group strings, or peptide-centric names. We therefore define
a canonical site key
\[
k = (\mathrm{protein\_accession}, \mathrm{residue}, \mathrm{position},
\mathrm{ptm\_type}),
\]
or, when the source does not fully localize a single residue, a canonical
site-group key
\[
k = (\mathrm{source\_site\_group\_id}, \mathrm{ptm\_type}).
\]

A reference alias table provides a deterministic mapping
\[
a(i) : \text{raw row } i \mapsto \text{canonical feature } k.
\]
If no alias table is supplied, the extractor falls back to direct construction
from the input columns. Ambiguous rows that map to multiple genes should be
handled explicitly by a policy such as \texttt{drop}, \texttt{split\_equal}, or
\texttt{first}; the recommended connectable default is \texttt{drop}.

\subsection{Extractor E1: \texttt{ptm\_site\_diff}}
\label{sec:ptm_site_diff}

\subsubsection{Base site statistic}
The extractor accepts several score definitions.

\paragraph{Custom score column.}
If the user already provides a site score column,
\[
x_i = c_i.
\]

\paragraph{Statistic mode.}
If a moderated $t$, $z$, or similar statistic is provided,
\[
x_i = t_i.
\]

\paragraph{Effect times significance.}
If a signed effect and a $p$-value are available,
\begin{equation}
x_i = d_i \cdot \min \left\{ -\log_{10}\!\bigl(\max(p_i,\varepsilon_p)\bigr),\,
C_p \right\},
\label{eq:ptm_effect_times_sig}
\end{equation}
where $\varepsilon_p > 0$ avoids taking the log of zero and $C_p$ caps
extremely small $p$-values.

\paragraph{Auto mode.}
In \texttt{auto} mode, prefer a supplied statistic, otherwise use
Equation~\eqref{eq:ptm_effect_times_sig} when both signed effect and
significance are available, otherwise fall back to a user-specified score
column.

\subsubsection{Confidence weighting}
Each site can be downweighted by confidence features:
\[
\omega_i = \omega_i^{(p)} \, \omega_i^{(\ell)} \, \omega_i^{(m)}.
\]

A simple implementation is:
\[
\omega_i^{(p)} =
\begin{cases}
\min\{-\log_{10}(\max(p_i,\varepsilon_p)), C_p\}/C_p, & p_i \text{ available},\\
1, & \text{otherwise},
\end{cases}
\]
\[
\omega_i^{(\ell)} =
\begin{cases}
\max\left\{0, \frac{\ell_i - \ell_{\min}}{1-\ell_{\min}}\right\}, &
\ell_i \text{ available},\\
1, & \text{otherwise},
\end{cases}
\]
\[
\omega_i^{(m)} =
\begin{cases}
\min\left\{1, \frac{\log(1+m_i)}{\log(1+m_{\max})}\right\}, &
m_i \text{ available},\\
1, & \text{otherwise}.
\end{cases}
\]

This weighting is optional. If the source table does not contain these fields,
set the corresponding factors to 1.

\subsubsection{Matched protein adjustment}
A major design choice is whether the signal should represent raw phosphosite
change or phosphosite change beyond total protein abundance change.

\paragraph{No adjustment.}
\[
\widetilde{x}_i = x_i.
\]

\paragraph{Direct subtraction.}
\begin{equation}
\widetilde{x}_i = x_i - \lambda_{\mathrm{prot}} \,
a_{g(a(i))}^{\mathrm{prot}},
\label{eq:ptm_protein_subtract}
\end{equation}
where $\lambda_{\mathrm{prot}}$ is typically 1 by default.

\paragraph{Residual mode.}
A more adaptive option estimates a global coefficient
\[
\widehat{\beta} = \arg\min_{\beta} \sum_i v_i
\left(x_i - \beta \, a_{g(a(i))}^{\mathrm{prot}}\right)^2,
\]
with optional weights $v_i$, then uses
\[
\widetilde{x}_i = x_i - \widehat{\beta} \,
a_{g(a(i))}^{\mathrm{prot}}.
\]

The recommended v1 default is the direct subtraction in
Equation~\eqref{eq:ptm_protein_subtract} because it is easy to explain,
deterministic, and usually enough to separate PTM-specific regulation from
protein-level abundance shifts.

\subsubsection{Reference ubiquity penalty}
An optional phosphosite reference bundle provides a detection-frequency prior.
Let $N_{\mathrm{ref}}$ be the number of reference samples used to build the
bundle and let $df_k$ be the number of those samples in which canonical site $k$
is detected. Define
\begin{equation}
u_k = \log \left( \frac{N_{\mathrm{ref}} + 1}{df_k + 1} \right).
\label{eq:ptm_idf}
\end{equation}

This is the phosphosite analogue of an inverse document frequency penalty:
ubiquitous sites receive smaller weight, rarer and therefore more selective
sites receive larger weight. When no bundle is available, set $u_k = 1$.

The v1 recommendation is to use only detection-frequency style priors and avoid
naively pooling raw phosphosite ratios across studies as if they were on a
common absolute scale.

\subsubsection{Signed or nonnegative site contribution}
Given a transformed score function $\psi$ selected by
\texttt{score\_transform}:
\[
\psi_{\mathrm{signed}}(x)=x,\quad
\psi_{\mathrm{abs}}(x)=|x|,\quad
\psi_{\mathrm{positive}}(x)=\max(x,0),\quad
\psi_{\mathrm{negative}}(x)=\max(-x,0),
\]
the final row-level contribution is
\begin{equation}
z_i = \psi(\widetilde{x}_i)\, \omega_i \, u_{a(i)}.
\label{eq:ptm_row_score}
\end{equation}

\subsubsection{Duplicate-row collapse to canonical sites}
Multiple peptide rows can map to the same canonical site. Collapse them using a
site-level policy:
\[
q_k = \operatorname{Agg}_{\mathrm{site}}
\{ z_i : a(i)=k \}.
\]

Useful options are:
\begin{itemize}
\item \texttt{highest\_confidence}: keep the row with largest
$\omega_i$ (or a deterministic proxy).
\item \texttt{max\_abs}: keep the row with largest $|z_i|$.
\item \texttt{mean} or \texttt{sum}: average or accumulate duplicate evidence.
\end{itemize}

The recommended default is \texttt{highest\_confidence} because many duplicate
rows are technical rather than biological repeats.

\subsubsection{Site-to-gene aggregation}
The core gene score is obtained from site-level scores $q_k$ assigned to gene
$g(k)$.

A naive sum,
\[
S_g^{\mathrm{sum}} = \sum_{k: g(k)=g} q_k,
\]
overweights proteins with many measured sites.

A naive mean,
\[
S_g^{\mathrm{mean}} = \frac{1}{n_g}\sum_{k: g(k)=g} q_k,
\]
can dilute a strong focal phosphosite.

A robust default is therefore a signed top-$K$ mean. Let
$T_g(K)$ be the subset of up to $K$ sites for gene $g$ with largest absolute
site score:
\[
T_g(K) \subseteq \{k : g(k)=g\}, \qquad
|T_g(K)| = \min\{K, n_g\}.
\]
Then define
\begin{equation}
S_g = \frac{1}{|T_g(K)|} \sum_{k \in T_g(K)} q_k.
\label{eq:ptm_signed_topk_mean}
\end{equation}

This keeps the sign information, limits site-count bias, and rewards coherent
multi-site regulation when several strong sites point in the same direction.

Recommended default:
\[
\texttt{gene\_aggregation} = \texttt{signed\_topk\_mean},
\qquad K = 3.
\]

\subsubsection{Gene set extraction}
Given gene scores $S_g$, select a compact program by one of the shared
repository strategies:
\begin{itemize}
\item \texttt{top\_k}: keep the genes with largest $|S_g|$.
\item \texttt{quantile}: keep genes above a chosen score quantile.
\item \texttt{threshold}: keep genes with $|S_g|$ above a minimum threshold.
\item \texttt{none}: retain all nonzero genes.
\end{itemize}

If a signed output is desired and the shared repository utilities support it,
apply within-set L1 normalization on the selected genes:
\begin{equation}
w_g = \frac{S_g}{\sum_{h \in \mathcal{G}_{\mathrm{sel}}} |S_h|}.
\label{eq:ptm_within_set_l1}
\end{equation}

The selected weighted set is written to \texttt{geneset.tsv}. When GMT export is
enabled with signed splitting, emit separate positive and negative sets:
\[
\mathcal{G}^+ = \{ g : S_g > 0 \},
\qquad
\mathcal{G}^- = \{ g : S_g < 0 \}.
\]

\subsubsection{Interpretation}
\begin{itemize}
\item Signed scores with no protein adjustment capture overall gain or loss of
measured PTM occupancy.
\item Signed scores with protein adjustment capture PTM-specific regulation
beyond total protein abundance changes.
\item Absolute or positive/negative transforms capture perturbation magnitude or
direction-specific evidence when downstream methods require nonnegative weights.
\item A ubiquity penalty shifts the extractor from ``many changed sites'' toward
``selective changed sites''.
\end{itemize}

\subsection{Extractor E2: \texttt{ptm\_site\_matrix}}
\label{sec:ptm_site_matrix}

Let $Y_{kj}$ be a site-by-sample abundance matrix, where $k$ indexes canonical
sites (or site-groups) and $j \in \{1,\ldots,n\}$ indexes samples. Let
$\gamma(j)$ denote a group label and let $\chi(j)$ denote a condition label.
The implemented matrix extractor computes within-study site-level statistics and
then reuses Section~\ref{sec:ptm_site_diff} onward.

\paragraph{Condition A vs condition B.}
For sample sets $A$ and $B$,
\[
d_k(A,B) =
\frac{1}{|A_k|}\sum_{j \in A_k} Y_{kj}
-
\frac{1}{|B_k|}\sum_{j \in B_k} Y_{kj},
\]
where $A_k$ and $B_k$ are the subsets of observed samples for site $k$ in the
two conditions.

A Welch-style statistic is
\[
t_k(A,B) =
\frac{d_k(A,B)}
{\sqrt{s_{k,A}^2/|A_k| + s_{k,B}^2/|B_k| + \varepsilon_t}}.
\]

Matched protein adjustment would use the analogous within-study contrast on a
matched protein matrix.

\paragraph{Grouped contrasts.}
The same site-level contrast estimation applies to:
\begin{itemize}
\item \texttt{group\_vs\_rest}: compare a focal sample group to all other
samples in the same matrix.
\item \texttt{condition\_within\_group}: estimate condition A vs B separately
inside each sample group and emit one program per group.
\item \texttt{baseline}: use the within-study site mean when no explicit
contrast is desired.
\end{itemize}

In all cases, the matrix front-end is only a contrast estimator. The
downstream canonicalization, confidence weighting, protein adjustment, ubiquity
penalty, duplicate collapse, and gene aggregation remain identical to
\texttt{ptm\_site\_diff}.

\subsection{Reference bundle design}
The recommended optional public bundle for the phosphoproteomics use case should
be intentionally small and derived from public site-level tables.

\subsubsection{Bundle contents}
\begin{itemize}
\item \texttt{phosphosite\_aliases\_human\_v1.tsv.gz}: maps raw feature labels to
canonical phosphosite keys and resolved genes.
\item \texttt{phosphosite\_ubiquity\_human\_v1.tsv.gz}: stores
$N_{\mathrm{ref}}$, $df_k$, fraction detected, and $u_k$ from
Equation~\eqref{eq:ptm_idf}.
\item \texttt{bundle\_provenance.json}: records public source studies and file
hashes used to build the compact tables.
\end{itemize}

\subsubsection{Why only a compact ubiquity bundle in v1}
For public phosphoproteomics, the safest cross-study prior is detection
frequency. Public relative-abundance tables can still be used inside a study for
contrasts, but v1 should avoid implying that raw ratios are directly comparable
across studies.

\subsection{Recommended defaults}
A practical default for phosphoproteomics is:
\begin{itemize}
\item \texttt{ptm\_type = phospho}
\item \texttt{score\_mode = auto}
\item \texttt{score\_transform = signed}
\item \texttt{protein\_adjustment = subtract}
\item \texttt{confidence\_weight\_mode = combined}
\item \texttt{site\_dup\_policy = highest\_confidence}
\item \texttt{gene\_aggregation = signed\_topk\_mean}
\item \texttt{gene\_topk\_sites = 3}
\item \texttt{select = top\_k}, \texttt{top\_k = 200}
\item \texttt{normalize = within\_set\_l1}
\item \texttt{emit\_gmt = true}, \texttt{gmt\_split\_signed = true}
\item optional bundle-backed ubiquity penalty when a local bundle is available
\end{itemize}

This default is intended to capture selective, direction-aware,
protein-adjusted phosphoregulation rather than raw site-count burden.

\subsection{CLI-to-theory crosswalk}
\begin{center}
\begin{tabular}{p{0.30\linewidth} p{0.28\linewidth} p{0.32\linewidth}}
\hline
CLI family & Mathematical object & Meaning \\
\hline
\texttt{--score\_mode}, \texttt{--score\_column}, \texttt{--logfc\_column},
\texttt{--padj\_column} & $x_i$ or $(d_i,t_i,p_i)$ &
Base site statistic \\
\texttt{--localization\_prob\_column}, \texttt{--peptide\_count\_column},
\texttt{--confidence\_weight\_mode} & $\omega_i$ &
Confidence weighting \\
\texttt{--protein\_logfc\_column}, \texttt{--protein\_adjustment} &
$a_g^{\mathrm{prot}}$, $\widetilde{x}_i$ &
Protein adjustment \\
\texttt{--study\_contrast}, \texttt{--group\_column},
\texttt{--condition\_column}, \texttt{--effect\_metric} &
$d_k(A,B), t_k(A,B)$ &
Within-study matrix contrast definition \\
\texttt{--resources\_manifest}, \texttt{--resources\_dir},
\texttt{--site\_ubiquity\_resource\_id} & $u_k$ &
Reference ubiquity prior \\
\texttt{--site\_dup\_policy} & $\operatorname{Agg}_{\mathrm{site}}$ &
Collapse duplicate rows to canonical sites \\
\texttt{--gene\_aggregation}, \texttt{--gene\_topk\_sites} &
$S_g$ &
Site-to-gene aggregation \\
\texttt{--select}, \texttt{--top\_k}, \texttt{--quantile},
\texttt{--min\_score}, \texttt{--normalize} & $w_g$ &
Compact program extraction and normalization \\
\hline
\end{tabular}
\end{center}

\subsection{Summary}
The implemented PTM extractor shares one conceptual target with the rest of the
repository: a compact, interpretable gene program. For PTM-omics, the most
useful public reference bundle in v1 is a compact alias-plus-ubiquity bundle
built from public phosphosite tables and used only for harmonization and
selectivity weighting.
